% Generated draft: requires mathematical and source review.
\documentclass[11pt]{article}
\usepackage[margin=25mm]{geometry}
\usepackage{fontspec}
\setmainfont{Latin Modern Roman}
\usepackage{amsmath,amssymb,mathtools,unicode-math}
\setmathfont{Latin Modern Math}
\usepackage{enumitem,needspace,xurl,hyperref}
\setlength{\parindent}{0pt}
\setlength{\parskip}{5pt}
\setlength{\emergencystretch}{3em}
\newcommand{\N}{\mathbb{N}}
\newcommand{\Z}{\mathbb{Z}}
\newcommand{\Q}{\mathbb{Q}}
\newcommand{\R}{\mathbb{R}}
\newcommand{\C}{\mathbb{C}}
\newcommand{\F}{\mathbb{F}}
\newcommand{\A}{\mathbb{A}}
\newcommand{\PP}{\mathbb{P}}
\newcommand{\E}{\mathbb{E}}
\newcommand{\eps}{\varepsilon}
\newcommand{\dd}{\,\mathrm{d}}
\begin{document}
\section*{UnsolvedMath ID: 30000789}
\textbf{Source problem code:} catalogue entry
\subsection{Definitions and mathematical statement}
For \(n\geq4\), let \(\mathcal R_n\) denote the space of algebraic curvature operators on \(\R^n\): self-adjoint operators on \(\Lambda^2\R^n\) satisfying the first Bianchi identity. Use the orthogonal decomposition
\[
R=R_I+R_{\operatorname{Ric}_0}+R_W,
\]
into scalar-curvature, trace-free-Ricci, and Weyl components. For \(c>0\), the cone in the restored source statement is
\[
\mathcal C_{n,c}=\{R\in\mathcal R_n:
\operatorname{scal}(R)\geq0,\quad
\|R_W\|^2\leq c\|R_I\|^2\}.
\]
No bound on \(R_{\operatorname{Ric}_0}\) is included in the supplied condition. Norms are understood as the usual induced operator-space norms; their precise source normalization and the convention for the quadratic operation \(R^\#\) should be confirmed. Forward invariance means that a solution of
\[
\frac{\dd R}{\dd t}=R^2+R^\#
\]
starting in \(\mathcal C_{n,c}\) remains in it at all later times in its interval of existence.

Let \(\mathcal T(n)\) mean that the B\"ohm--Wilking classification referred to in the report gives the exact set of \(c>0\) for which \(\mathcal C_{n,c}\) is invariant. The question concerns the least integer threshold
\[
n_0=\min\{N:\mathcal T(n)\text{ holds for every integer }n\geq N\}.
\]
The source reports \(n_0\geq12\) and asks whether \(n_0=12\). This is an eventual threshold over all larger dimensions, not only a question in dimension twelve. The supplied record does not reproduce the classification theorem or its proposed set of admissible \(c\). Thus \(\mathcal T(n)\), and hence the threshold, remains a reference-dependent predicate until that missing theorem is supplied; the existence of a finite threshold is part of the reported theorem context, not established in this draft.
\subsection{Short English statement}
Is twelve the sharp dimension from which the cited B\"ohm--Wilking classification gives precisely the invariant cones defined by nonnegative scalar curvature and \(\|R_W\|^2\leq c\|R_I\|^2\)? The exact classification theorem is missing from the extraction.
\subsection{Sources}
\begin{itemize}
\item Local UnsolvedMath registry, permanent ID \(30000789\), source code catalogue entry, ``Sharp Dimension Threshold for Ricci-Flow Curvature Conditions.'' The cleaned statement attributes the classification to B\"ohm--Wilking; no individual proposer is supplied.
\item \emph{Differentialgeometrie im Grossen} (2008), Oberwolfach Reports, DOI 10.4171/owr/2007/32, pp. 1865--1912: \url{https://doi.org/10.4171/owr/2007/32}. This is the sole supplied external reference.
\item The literature triage checked 2026-08-21 reports the lower bound n\_0 at least 12 and no located authoritative proof of equality or counterexample. Neither the bound nor present-day status is certified by this draft.
\end{itemize}
\end{document}
